%-------------------------------------------------------------------------------
%	SECTION TITLE
%-------------------------------------------------------------------------------
\cvsection{个人项目}
%-------------------------------------------------------------------------------
%	CONTENT
%-------------------------------------------------------------------------------
\begin{cventries}
  \cventry
  {University of Toronto}
  {Dinosaur Game on STM32 MCU} % Name
  {C, STM32, Embedded Systems}
  {2026年4月} % Date(s)
  {
    \begin{cvitems} % Description(s)
      \item{使用C语言在STM32F446ZE MCU上开发了一款2D Dinosaur Game, 实现了游戏玩法、GPIO和I2C外设驱动程序, 以及双缓冲视频引擎, 并在3周内交付了一款可运行的游戏。}
      \item{集成了一个 4$\times$3 Keypad, 用于实时用户控制, 并利用CPU中断实现低延迟的输入处理。}
      \item{开发了一个轻量级、无OS的实时渲染2D图形引擎, 支持同步驱动2块SSD1306 OLED显示屏, 使用了Double Buffering和DMA技术实现了基于I2C接口的低延迟高稳定性的实施画面传输。}
    \end{cvitems}
  }

  \cventry
  {University of Toronto}
  {Runner Game (FPGA Board Game)} % Name
  {C, RISC-V Assembly, CPUlator}
  {2025年3月} % Date(s)
  {
    \begin{cvitems} % Description(s)
      \item{在DE1-SoC FPGA Board上使用Verilog开发了一个2D Runner Game, 实现核心游戏逻辑、PS/2 Keyboard、音频组件和VGA组件, 并在三周内交付了可玩的游戏。}
      \item{集成了PS/2 Keyboard以实现实时用户控制, 利用CPU Interupt实现低延迟输入处理。}
      \item{构建了一个轻量、无OS的2D引擎, 以60帧每秒渲染彩色图像、动画和文字。}
      \item{开发了一个基本的Square Wave Synthesizer, 在游戏事件(如得分和完成游戏)中实现动态音效。}
    \end{cvitems}
  }

  \cventry
  {University of Toronto}
  {Greedy Mouse Game (FPGA Board Game)} % Name
  {Verilog, FPGA Board, ModelSim}
  {2024年11月} % Date(s)
  {
    \begin{cvitems} % Description(s)
      \item{在DE1-SoC FPGA Board上用Verilog设计并实现了一个2D Greedy Mouse Game, 具备PS/2 Keyboard输入和VGA图形输出功能。}
      \item{开发了一个PS/2 Keyboard Controller以实现实时用户输入, 并将其与FSM和片上存储模块集成, 用于驱动VGA输出, 实现位图动画的流畅显示并增强游戏视觉效果。}
      \item{构建了一个轻量、无OS的2D引擎, 以60帧每秒渲染彩色图像、动画和文字。}
      \item{使用ModelSim和DESim进行早期阶段的仿真和调试, 最终使用Quartus Prime完成综合并部署到FPGA Board, 并进行了反复优化。}
    \end{cvitems}
  }

  \cventry
  {University of Toronto}
  {GIS Route Optimization Application} % Name
  {C++, GTK, Git, A*, Dijkstra}
  {2025年1月 - 2025年4月} % Date(s)
  {
    \begin{cvitems} % Description(s)
      \item{在课程项目的三人小组中, 使用C++和GTK在Mate桌面环境下开发了一个GIS(地图)桌面应用, 实现了地图渲染、地名搜索、最短路径和多停点路径查找功能。}
      \item{利用Multithreading预处理4GB原始坐标数据, 将其转换为结构化的点、线和多边形格式, 以便在16核测试机器上在50秒内完成从数据解析到画布渲染的全部操作。}
      \item{实现了包括A*的最短路径计算算法, 以及Dijkstra、Multi-start Greedy和Simulated Annealing, 获得了课程技术分数的90\%。}
      \item{通过使用TomTom API并整合libcurl的HTTP模块, 实现了实时交通可视化。}
    \end{cvitems}
  }

  \cventry
  {GitHub 开源项目}
  {贸易流管理系统} % Name
  {TypeScript, Node.js, PostgreSQL, React}
  {2025年7月 - 现在} % Date(s)
  {
    \begin{cvitems} % Description(s)
      \item{构建了采用模块化单仓库结构的全栈式交易管理解决方案, 具备基于JWT的RBAC认证机制及 i18n 支持, 以实现订单与财务管理的无缝衔接。}
      \item{采用Node.js和Express构建高性能后端系统, 运用Decimal.js实现精准金融计算, 并借助ESBuild优化构建产物。}
      \item{采用React和Ant Design开发响应式前端, 采用组件驱动架构以加速功能交付并保持代码复用性}
    \end{cvitems}
  }

  \cventry
  {University of Toronto}
  {LPR 系统} % Name
  {Python, PyTorch}
  {2026年1月 - 2026年4月} % Date(s)
  {
    \begin{cvitems} % Description(s)
      \item{构建了一个基于深度学习的车牌识别系统, 设计了CNN+CRNN神经网络, 编写了用于训练和推理的PyTorch代码, 并完成了模型的训练和测试。}
      \item{构建了一个Python管道, 用于批量生成车牌图像、为其添加噪点并将其与背景拼接, 自动生成10,000个包含完整车辆背景的训练数据点。}
      \item{利用Python和uv作为包管理器, 建立了一套全面的数据生成和训练管道, 并编写了完整的项目文档, 以确保训练结果具有高度可重复性。}
    \end{cvitems}
  }
\end{cventries}
